\newpage
\pdfbookmark[1]{\infoname}{info} % Lisame infolehe PDF-faili järjehoidjatesse.

% Eesti keeles info.
\begin{info}
\begin{abstract}
Eestikeelse informaatika lõputöö kirjutamine on tudengile keeleliselt ja struktuurselt nõudlik ülesanne, sest valdkonna terminoloogia on inglise keeles välja kujunenud ning akadeemilise stiili reegleid kohaldatakse eesti keeles erinevalt. Käesoleva bakalaureusetöö eesmärk on luua ja hinnata suurel keelemudelil põhinev prototüüp, mis annab eestikeelsele informaatika lõputöö tekstile automaatset tagasisidet neljas kategoorias: struktuur, akadeemiline stiil, terminoloogia järjepidevus ja viitamisvajadus. Töös käsitletakse suurte keelemudelite võimekust eestikeelse tehnilise teksti analüüsimisel ning struktureeritud promptimise mõju tagasiside kvaliteedile. Prototüüp on realiseeritud veebirakendusena, mis võtab sisendiks tudengi kirjutatud peatüki, suunab selle valitud keelemudelile struktureeritud viibaga ning kuvab kasutajale kategoriseeritud tagasiside. Süsteemi hindamiseks koostati 20 tekstikatkendist koosnev testkogu, milles probleemid märgendati käsitsi koos juhendajaga. Mudeli väljundit võrreldi kuldstandardiga täpsuse, saagise ja F\textsubscript{1}-skoori abil. Tulemused näitavad, et struktureeritud promptimine annab kõikides kategooriates üldisest viibast täpsemat tagasisidet, kusjuures mudel on kõige usaldusväärsem viitamisvajaduse ja struktuuri tuvastamisel ning kõige problemaatilisem terminoloogia järjepidevuse hindamisel. Töös rõhutatakse, et süsteem on tudengi enda kirjutatud teksti tagasisidevahend, mitte juhendaja ega tekstigeneraator.
\end{abstract}

\keywords{}

\cercs{P175 Informaatika, süsteemiteooria}
% Koodid leiate: https://www.etis.ee/Portal/Classifiers/Index/26
\end{info}


% Inglise keeles info.
\begin{otherInfo}{english}{Using Large Language Models for Quality Control of Estonian Technical Computer Science Texts}
\begin{abstract}
Writing a computer science thesis in Estonian is linguistically and structurally demanding for a student because the field's terminology has evolved in English and academic style conventions are applied differently in Estonian. The aim of this bachelor's thesis is to design and evaluate a prototype based on a large language model (LLM) that provides automatic feedback on Estonian-language computer science thesis texts in four categories: structure, academic style, terminological consistency, and the need for citation. The thesis examines the capability of large language models to analyse Estonian technical text and the effect of structured prompting on the quality of the resulting feedback. The prototype is implemented as a web application that takes a student-written chapter as input, sends it to the chosen language model with a structured prompt, and presents categorised feedback to the user. To evaluate the system, a test set of 20 text excerpts was created, with issues annotated manually together with the supervisor. The model's output was compared against this gold standard using precision, recall, and F\textsubscript{1}-score. The results indicate that structured prompting yields more accurate feedback than a generic prompt across all categories, that the model is most reliable when detecting the need for citation and structural problems, and least reliable when assessing terminological consistency. The thesis emphasises that the system is a feedback tool for text the student has written themselves; it neither replaces the supervisor nor generates the thesis on the student's behalf.
\end{abstract}

\keywords{Large language models, Estonian language, academic text, feedback, quality control, prompting}

\cercs{P175 Informatics, systems theory}
\end{otherInfo}
